% =====================================================================
%  PLANTILLA · documento article (matemática) — suite LaTeX
%  Copiá este archivo y empezá a escribir donde dice "CONTENIDO".
%  Compilar:  latexmk -pdf article.tex
%  Base estándar comentada (ver doc "Estándares y buenas prácticas").
% =====================================================================
\documentclass[11pt,a4paper]{article}

% --- Codificación y fuentes (motor pdfLaTeX) ------------------------
\usepackage[utf8]{inputenc}   % entrada UTF-8 (acentos en el .tex)
\usepackage[T1]{fontenc}      % salida T1: guionado y copy-paste correctos
\usepackage{lmodern}          % fuentes vectoriales (no bitmaps)
\usepackage{microtype}        % pulido tipográfico (protrusion + expansion)

% --- Idioma ---------------------------------------------------------
\usepackage[spanish,es-noshorthands]{babel} % es-noshorthands: no rompe tikz/urls
\usepackage{csquotes}                        % \enquote{...}: comillas correctas

% --- Matemática -----------------------------------------------------
\usepackage{mathtools}  % superconjunto de amsmath (+ \coloneqq, delimitadores)
\usepackage{amssymb}    % \mathbb, \mathcal, relaciones, flechas
\usepackage{amsthm}     % entornos de teorema/definición/demostración
\usepackage{bm}         % negrita matemática correcta: \bm{x}
\usepackage{siunitx}    % números y unidades: \num, \qty

% --- Tablas, gráficos, layout ---------------------------------------
\usepackage{booktabs}\usepackage{tabularx}   % tablas limpias
\usepackage{graphicx}                         % \includegraphics
\usepackage{pgfplots}\pgfplotsset{compat=1.18}% gráficos de funciones/datos
\usepackage{caption}\usepackage{subcaption}\usepackage{float}
\usepackage[margin=2.5cm]{geometry}           % márgenes
\usepackage{enumitem}                         % listas personalizables

% --- Bibliografía (descomentar si se usa) ---------------------------
% \usepackage[backend=biber,style=numeric]{biblatex}
% \addbibresource{refs.bib}

% --- Enlaces y referencias (casi al final) --------------------------
\usepackage{hyperref}
\hypersetup{unicode, hidelinks,
            pdftitle={Título del documento}, pdfauthor={Tu Nombre}}
\usepackage[spanish,capitalise,nameinlink]{cleveref} % DESPUÉS de hyperref

% --- Macros propias (DRY + semántica) -------------------------------
\newcommand{\R}{\mathbb{R}}  \newcommand{\N}{\mathbb{N}}
\newcommand{\Z}{\mathbb{Z}}  \newcommand{\Q}{\mathbb{Q}}
\DeclareMathOperator{\sen}{sen}                 % seno (español, redonda)
\DeclarePairedDelimiter{\abs}{\lvert}{\rvert}   % \abs{x}, \abs*{\frac..}
\DeclarePairedDelimiter{\norm}{\lVert}{\rVert}
\theoremstyle{plain}      \newtheorem{teo}{Teorema}[section]
\theoremstyle{definition} \newtheorem{defi}[teo]{Definición}

\title{Título del documento}
\author{Tu Nombre}
\date{\today}

\begin{document}
\maketitle
% \tableofcontents \newpage    % índice en hoja aparte, si lo necesitás

% =====================================================================
% CONTENIDO  — borrá los ejemplos y escribí lo tuyo
% =====================================================================
\section{Introducción}
Texto con una cita \enquote{entre comillas} y una referencia a
\cref{sec:resultados}.

\section{Resultados}\label{sec:resultados}
Matemática en línea $f(x)=x^2$ y en bloque:
\begin{equation}\label{eq:cuad}
  f(x) = ax^2 + bx + c, \qquad a \neq 0 .
\end{equation}
Por \cref{eq:cuad}, el discriminante es $\Delta = b^2 - 4ac$.

\begin{teo}\label{teo:ej}
  Para todo $x \in \R$ se cumple $\abs{x} \ge 0$.
\end{teo}
\begin{proof}
  Inmediato por la definición de valor absoluto.
\end{proof}

\begin{figure}[H]
  \centering
  \begin{tikzpicture}
    \begin{axis}[width=8cm, height=5cm, axis lines=middle,
                 xlabel=$x$, ylabel=$y$, samples=80]
      \addplot[blue, thick, domain=-2:2] {x^2};
    \end{axis}
  \end{tikzpicture}
  \caption{La parábola $y=x^2$.}\label{fig:par}
\end{figure}

\begin{table}[H]
  \centering
  \begin{tabular}{@{}lS@{}}
    \toprule
    Magnitud & {Valor} \\
    \midrule
    Ancho  & 12.5 \\
    Alto   &  3.2 \\
    \bottomrule
  \end{tabular}
  \caption{Una tabla con \texttt{booktabs}.}\label{tab:ej}
\end{table}

% \printbibliography   % si usás biblatex

\end{document}
